\section{记忆管理：遗忘策略}

真正的记忆系统不只是不断积累，还需要\textbf{精心筛选}。一个不断增长、缺乏重点的存储会随时间退化——检索噪声增大、延迟攀升、矛盾的记忆让 Agent 困惑。

\begin{warningbox}{无限堆积记忆的陷阱}
没有遗忘策略的记忆系统最终会变得\textbf{比没有记忆更糟糕}：检索噪声淹没有效信息，互相矛盾的记忆导致 Agent 决策混乱。
\end{warningbox}

\subsection{Time-based Decay（基于时间的衰减）}

较旧的记忆通常相关性较低。通过 recency 和 semantic relevance 的组合来评分记忆。研究中常用的公式：

$$
\text{score} = \alpha \cdot \text{relevance} + (1 - \alpha) \cdot \text{recency\_decay}(t)
$$

其中：
\begin{itemize}[leftmargin=2em]
    \item $\alpha$：语义相关性与时间衰减之间的权衡系数
    \item $\text{relevance}$：查询与记忆之间的语义相似度
    \item $\text{recency\_decay}(t)$：基于时间 $t$ 的衰减函数（如指数衰减）
\end{itemize}

\subsection{Importance Scoring（写入时重要性评分）}

存储记忆时，让模型对自己的输出进行重要性评分，只存储高分项。这从源头过滤噪声。

\subsection{Periodic Consolidation（定期合并）}

运行定期任务（如每天夜间），将重复或高度相似的记忆合并为一条规范的摘要。这类似于人类睡眠过程中的\textbf{记忆巩固}机制。

\begin{knowledgebox}{类比人类记忆}
人类大脑在睡眠中进行"记忆巩固"——将白天零散的短期记忆整理、合并为长期记忆。Agent 的 Periodic Consolidation 机制借鉴了同样的原理：定期清理冗余，保留精华。
\end{knowledgebox}

\subsection{本章小结}

三种遗忘策略各有用途：Time-based Decay 让旧记忆自然淡出，Importance Scoring 在写入时过滤噪声，Periodic Consolidation 定期合并重复记忆。三者结合可维持高质量的记忆库。

